\section{Notation and setup}
\label{sec:prelim}

We collect the notation, the formal setup, and the standard
arithmetic conventions used throughout the rest of the paper.

\paragraph{Notation.}
We work in $\R^d$ with the Euclidean norm $\norm{\cdot}$ and the inner
product $\inner{\cdot}{\cdot}$. Throughout, $c, C, c_1, c_2, \ldots$
denote universal positive constants whose values may change from line to
line and depend only on quantities specified in each lemma's hypothesis.
Constants that depend on additional problem parameters carry those
parameters as subscripts (e.g.\ $C_{B_U, M}$).

We fix a finite set $F$ of \emph{questions} (the ``question field''), a
single transformer layer index $i$ and a single attention-head index $h$
within that layer; these are fixed throughout the paper. The argument
extends to any other $(i, h)$ pair by union bound at logarithmic cost in
the total number of (layer, head) pairs (see \Cref{rem:per_layer_extension}
in \Cref{sec:main_theorem}).

\paragraph{The reasoning trajectory.}
For a question $Q \in F$, the reasoning model generates a sequence of
tokens between the markers $\langle\mathrm{think}\rangle$ and
$\langle/\mathrm{think}\rangle$. We index the generated tokens by
$j = 1, 2, \ldots$ in emission order. We write $\mathcal F_{j-1}$ for the
$\sigma$-algebra generated by the model's stochastic sampling decisions up
to (but not including) step $j$; the sampling at step $j$ is then
$\mathcal F_{j-1}$-conditional.

For each $j \ge 1$, let $k_j \in \R^{d_k}$ and $V_j \in \R^d$ denote the
\emph{key} and \emph{value} vectors produced by layer $i$, head $h$ at the
$j$-th reasoning position, and let $q \in \R^{d_k}$ denote the query
vector at the $\langle/\mathrm{think}\rangle$ position (the position at
which the model commits to an answer). The key, value, and query vectors
depend on the model parameters (fixed) and on $Q$ and the reasoning
trajectory generated so far.

\begin{definition}[Cumulative softmax attention]\label{def:softmax_attention}
The single-head softmax attention output of layer $i$, head $h$, at the
$\langle/\mathrm{think}\rangle$ position, restricted to the first $j$
reasoning keys, is
\begin{align}\label{eq:cumulative_softmax}
x_j \;\coloneqq\; \sum_{k=1}^{j} w_{j,k} V_k,
\qquad
w_{j,k} \;\coloneqq\; \frac{e^{\inner{q}{k_k}}}{s_j},
\qquad
s_j \;\coloneqq\; \sum_{i=1}^{j} e^{\inner{q}{k_i}},
\end{align}
where $\inner{q}{k_k}$ denotes the rescaled inner product $q^\top k_k / \sqrt{d_k}$
absorbing the standard $\sqrt{d_k}$ normalisation factor of
\cite{vaswani2017attention}. We write the per-step contribution as
\begin{align}\label{eq:gj_def}
   g_j \;\coloneqq\; \frac{e^{\inner{q}{k_j}}}{s_j}\, V_j,
\end{align}
so that the recurrence $x_j = (s_{j-1}/s_j)\, x_{j-1} + g_j$ holds (see
\Cref{lem:softmax_running_average} for the formal identity).
\end{definition}

\begin{remark}\label{rem:notation_recurrence}
With $x_0 \coloneqq 0$, the recurrence form
$x_j = (s_{j-1}/s_j)\, x_{j-1} + g_j$
is equivalent to \Cref{def:softmax_attention}. This is the
shape that motivates the ``one-step (biased) gradient update'' framing
of the $x_j$-dynamics formalised by the assumption and analysis in the
following sections.
\end{remark}

\paragraph{Vocabulary, unembedding, and the constrained-softmax loss.}
Let $\mathcal V$ denote the model's tokenizer vocabulary, a fixed finite
set of tokens. Let $W_U \in \R^{|\mathcal V| \times d}$ denote the
(fixed) unembedding matrix that maps a $d$-dimensional latent vector to
logits over the tokenizer vocabulary. For each $v \in \mathcal V$ we
write $W_U^{(v)} \in \R^d$ for the $v$-th row of $W_U$ (viewed as a
column vector).
For each $x \in \R^d$, define
the softmax logits and posterior over the vocabulary:
\begin{align}\label{eq:z_p_def}
   z(x) \;\coloneqq\; W_U\, x \;\in\; \R^{|\mathcal V|},
   \qquad
   p(x) \;\coloneqq\; \softmax(z(x)) \;\in\; \Delta^{|\mathcal V|},
\end{align}
where $\Delta^{|\mathcal V|}$ is the probability simplex over $\mathcal V$.

\begin{definition}[Correct token set]\label{def:correct_set}
For each question $Q \in F$, the \emph{correct-token set} is a non-empty
subset $\Cset(Q) \coloneqq \Correct(Q) \subseteq \mathcal V$ of vocabulary
tokens that decode as correct answers for $Q$. We write
$\1_{\Cset(Q)} \in \{0,1\}^{|\mathcal V|}$ for its indicator vector.
\end{definition}

\begin{definition}[Correct-mass and constrained-softmax loss]\label{def:loss}
For each $Q \in F$ and each $x \in \R^d$, define the
\emph{correct-token softmax mass} and the
\emph{constrained-softmax loss}:
\begin{align}\label{eq:loss_def}
   \cmass(x; Q) \;\coloneqq\; \sum_{v \in \Cset(Q)} p(x)_v
                          \;=\; \1_{\Cset(Q)}^{\top}\, p(x),
   \qquad
   \loss(x; Q) \;\coloneqq\; -\log \cmass(x; Q).
\end{align}
The $\Cset(Q)$-conditional posterior over $\mathcal V$ is
\begin{align}\label{eq:qcond_def}
   \qcond(x; Q)_v \;\coloneqq\;
   \frac{p(x)_v \cdot \1[v \in \Cset(Q)]}{\cmass(x; Q)},
   \qquad v \in \mathcal V,
\end{align}
so $\qcond(x; Q) \in \Delta^{|\mathcal V|}$ is supported on $\Cset(Q)$
(it is the posterior $p(x)$ conditioned on the event $\{v \in \Cset(Q)\}$).
Both $p$ and $\qcond$ are probability vectors on $\mathcal V$.
\end{definition}

\begin{remark}\label{rem:loss_nondegenerate}
$\cmass(x; Q) > 0$ for every $x \in \R^d$ since $p(x)$ is the softmax of
a finite vector and $\Cset(Q)$ is non-empty; consequently
$\loss(x; Q) < \infty$ for every $x$, and $\loss(x; Q) \ge 0$ since
$\cmass(x; Q) \le 1$. The minimum value $\loss = 0$ corresponds to the
correct set carrying all the softmax mass, $\cmass = 1$.
\end{remark}

\paragraph{Decoding.}
The model's decoding map applied to a latent vector $x \in \R^d$ returns
a final answer token
\begin{align}\label{eq:decode_def}
   \dec(x) \;\coloneqq\; \arg\max_{v \in \mathcal V} (W_U x)_v,
\end{align}
breaking ties arbitrarily (e.g.\ by vocabulary index). The decode is
\emph{correct} for $Q$ if $\dec(x) \in \Cset(Q)$.

\begin{remark}\label{rem:decode_loss_relationship}
The decoding correctness event $\{\dec(x) \in \Cset(Q)\}$ is a discrete
event on the polyhedral cone
$\{x : \max_{v \in \Cset(Q)} (W_U x)_v \ge \max_{v \notin \Cset(Q)} (W_U x)_v\}$,
while the loss $\loss(x; Q)$ is a smooth function of $x$. We connect the
two in \Cref{lem:logit_margin_decoding}: a sufficient (not necessary)
condition for correct decoding is that $\loss(x; Q)$ is small enough.
This is the bridge from the optimisation analysis of $\loss$ to the
decoding guarantee that appears in the headline theorem.
\end{remark}
